\begin{Heading}{chapter}
  \Title{Quick Start}
  \RefLabel{chap:quickstart}
\end{Heading}

Thank you for using \CoCoTeX.

In this introductory chapter, we are getting familiar with \CoCoTeX's
design principles and how to use it to define the following boxes:

\begin{Info}[\Title{Info box}]
  This box holds some general information.
\end{Info}

Source code:
\begin{lstlisting}
\begin{Info}[\Title{Info box}]
  This box holds some general information.
\end{Info}
\end{lstlisting}

\begin{Background}
  This box type holds background information if you are interested
  into the inner workings of the {\CoCoTeX} framework. It can also
  have a title.
\end{Background}

Source code:
\begin{lstlisting}
\begin{Background}
  This box type holds background information if you are interested
  into the inner workings of the {\CoCoTeX} framework. It can also
  have a title.
\end{Background}
\end{lstlisting}

\begin{Alert}[\Title{Alert box}]
  This box holds important information and common error sources.
\end{Alert}

Source code:
\begin{lstlisting}
\begin{Alert}[\Title{Alert box}]
  This box holds important information and common error sources.
\end{Alert}
\end{lstlisting}

\begin{Alert}[\Icon{\faCircleExclamation}]
  Alert box with an alternative icon but no title
\end{Alert}

Source code:
\begin{lstlisting}
\begin{Alert}[\Icon{\faCircleExclamation}]
  Alert box with an alternative icon no title
\end{Alert}
\end{lstlisting}
As you can see, each of those boxes contain three pieces of
information: They may or may not have a \textbf{title}, they have an
\textbf{icon} that can be set freely, and they have a mandatory
\textbf{content}.


The \textit{general idea} of a “box” is what in the {\CoCoTeX}
Framework is called a \textbf{Container}. Each box you can see in this
pdf is an \textbf{Instance} of that Container. Each of those
Containers is defined by the three informational data pieces they
share, namely the \textit{Title}, the \textit{Icon} and the
\textit{Content}. We call those pieces of Information
\textbf{Components} of a Container. The instructions, how those
Components are typeset is stored in so-called
\textbf{Properties}. Both Components and Properties are called the
\textbf{Types} of a Container.

In this documentation, there are three “flavours“ of the more abstract
Container “Box”, namely \textit{Info}, \textit{Background}, and
\textit{Alert}. We call those flavours \textbf{Subcontainer} or
\textbf{Child Containers}, since they \textbf{inherit} the Properties
and Components from the abstract Box Container.

We start with the most abstract Container \textbf{Box}. TO define a
new Container, we use the \UsageMacro{\ccDeclareContainer} macro which
takes two arguments, the \textit{name} of the Container and the
\textit{body} of the declaration.
\begin{lstlisting}
\ccDeclareContainer{Box}{%
  …
}
\end{lstlisting}
In the body, we declare the two Types that define this Container using
the \UsageMacro{\ccDeclareType} macro, which also takes two argumets, the name of the Data Type and its declaration body:
\begin{lstlisting}
\ccDeclareContainer{Box}{%
  \ccDeclareType{Components}{…}
  \ccDeclareType{Properties}{…}
}
\end{lstlisting}
We start with the Components. Each Component is declared with the
\UsageMacro{\ccDeclareComponent} macro. This macro has three mandatory
argument, althouggh in this tutorial we use only the first one which
holds the \textit{name} of the Component.
\begin{lstlisting}
\ccDeclareContainer{Box}{%
  \ccDeclareType{Components}{%
    \ccDeclareComponent{Title}{}{}
    \ccDeclareComponent{Icon}{}{}
    \ccDeclareComponent{Content}{}{}
  }
  \ccDeclareType{Properties}{…}
}
\end{lstlisting}
As mentioned before, the Properties define how the Components are
processed. Let us think about how we could to process our components:
Maybe, we want to manipulate the size of the icon and its color. So,
let us invent Properties that hold those kinds of information. To
define a Property, we use the macro \UsageMacro{\ccSetProperty} which
takes two arguments, first name of the Property and second its
definition. As a convention, Property names are lower-case and words
are separated by hyphen:
\begin{lstlisting}
\ccDeclareContainer{Box}{%
  \ccDeclareType{Components}{…}
  \ccDeclareType{Properties}{
    \ccSetProperty{icon-color}{black}
    \ccSetProperty{icon-size}{\fontsize{24bp}{32bp}}
  }
}
\end{lstlisting}
Finally, we need to define how the Icons are printed. We want the
icons to overlap the box on the left-side and we don't want the icon
to spread the lines. for that, we use the plain\TeX primitives
\lstinline{\llap} and \lstinline{\smash}. The icon itself is placed in
a narrow vbox that grows into its depth (\lstinline{\vtop}) that we
need to lower a bit. Inside that vbox, we print the Icon. Remember
that the Icon is a Component. To print the content of a component, we
use the \UsageMacro{\ccUseComp} macro which take a Component's name as
argument.
\begin{lstlisting}
\ccDeclareContainer{Box}{%
  \ccDeclareType{Components}{…}
  \ccDeclareType{Properties}{%
    …
    \ccSetProperty{icon-format}{%
      \llap{\smash{\lower1ex\vtop{%
            \hsize32bp
            \ccUseProperty{icon-size}\selectfont%
            \color{\ccUseProperty{icon-color}}%
            \ccUseComp{Icon}%
          }% /vtop
        }% /smash
      }% /llap
    }% /ccSetProperty{icon-format}
  }% /ccDeclareType{Properties}
}
\end{lstlisting}
As you can see, properties are small snippets of {\LaTeX} code that
sort, format and use Components, but also other Properties.

We also need some properties to process the Title Component. We want
the Title to set in bold-face and sans-serif font. For this, we
declare a property \texttt{title-face}. Another property,
\texttt{title-format}, combines the \texttt{font-face} Property with
the Title Component:
\begin{lstlisting}
\ccDeclareContainer{Box}{%
  \ccDeclareType{Components}{…}
  \ccDeclareType{Properties}{%
    \ccSetProperty{title-face}{\bfseries\sffamily}
    \ccSetProperty{title-format}{%
      \bgroup
        \ccUseProperty{title-face}
        \ccUseComp{Title}
      \egroup
    }%
  }% /ccDeclareType
}% /ccDeclareContainer
\end{lstlisting}
In order to make a Container instanciable, we need to declare a
corresponding {\LaTeX} environment using the \lstinline{\ccDeclareEnv}
macro. This macro takes two arguments, the meaning of which will be
explained in a bit.

\begin{Background}
  Strictly speaking, \texttt{\string\ccDeclareEnv} does not
  define a {\LaTeX} environment \texttt{<container\_name>}, but two
  macros \protect\texttt{\textbackslash<container\_name>} and
  \texttt{\textbackslash end<container\_name>}.
\end{Background}

By default, the environment is conveniently called the same as the
Container:
\begin{lstlisting}
\ccDeclareContainer{Box}{%
  \ccDeclareType{Components}{…}
  \ccDeclareType{Properties}{…}
  \ccDeclareEnv{}{}
}
\end{lstlisting}
The environment declaration alone does nothing on its own, yet, but we
will look into that later.

Let us first declare out Sub Containers, starting with the Background
Container:
\begin{lstlisting}
\ccDeclareContainer{Background}{%
}
\end{lstlisting}
Now, we need to tell the Background Container that it should inherit
the Components and Properties from the Box Container. For that, we use
the \UsageMacro{\ccInherit} macro, which takes two arguments. First
the Types we want to inherit, and second the Container from which we
want to inherit the Types.
\begin{lstlisting}
\ccDeclareContainer{Background}{%
  \ccInherit{Components,Properties}{Box}
}
\end{lstlisting}
Now, we want to specify two things that are prototypical for the
Background boxes: They have the Atom symbol as icon, and their color
is always green. We can do this by utilizing the Parent Container's
Types:
\begin{lstlisting}
\ccDeclareContainer{Background}{%
  \ccInherit{Components,Properties}{Box}
  \ccDeclareType{Components}{\Icon{\faAtom}}
  \ccDeclareType{Properties}{
    \ccSetProperty{icon-color}{green}
  }
}
\end{lstlisting}
Since we already inherited the Box's Components, we can immediately
use the Icon Component and give it a fix value.

By re-declaring an already existing Property name
(\texttt{icon-color}), we can override the value that came from the
Parent Container.

The only thing left to do is to declare Background's environment. For
that, we now use the two arguments of the \UsageMacro{\ccDeclareEnv}
macro. The first argument tells {\CoCoTeX} what should happen at the
begin, and the second what should happen before the end of the
environment:
\begin{lstlisting}
\ccDeclareContainer{Background}{%
  \ccInherit{Components,Properties}{Box}
  \ccDeclareType{Components}{…}
  \ccDeclareType{Properties}{…}
  \ccDeclareEnv{\begin{Box}}{\end{Box}}
}
\end{lstlisting}
With that, we define an environment \texttt{Background} that does
nothing on its own but to call the parent container's own environment.

We can now start to process the parent's environment. Since
\lstinline{\ccDeclareEnv} already defines the low-level macros that
constitute a {\LaTeX} enviroment, we need to use
\lstinline{\RenewDocumentEnvironment} (or
\lstinline{\DeclareDocumentEnvironment}).

At the beginning of that environment, we need to invoke the macro
\lstinline{\ccEvalType} which takes one argument, the name of a
Type. We need to call this macro twice, one time for each of your two
data types \texttt{Properties} and \texttt{Components}.
\begin{lstlisting}
\RenewDocumentEnvironment{Box}{ }{
  \ccEvalType{Properties}%
  \ccEvalType{Components}%
}{%
}
\end{lstlisting}
TODO: explain briefly the concept behind Type evaluation.
